\chapter{Free Particles, Fourier Evolution, and the Classical Limit}

\section{When Classical Repair Was No Longer Enough}

There are two ways to respond when an old theory fails. One can try to repair
the mechanism, adding enough machinery to preserve the familiar picture, or
one can search for a principle that changes the picture itself. The
Michelson--Morley experiment encouraged both reactions. Elaborate mechanical
models of the ether could be made to imitate the observed result, but
Einstein's relativity began from a sharper principle and made the old
mechanism unnecessary. Quantum theory followed a similarly reluctant path
from repair to reconstruction.

For the present discussion, the decisive classical failure is blackbody
radiation. Electromagnetic radiation in a cavity can be decomposed into
normal modes. Each mode behaves mathematically like a harmonic oscillator,
and classical equipartition assigns to it an average thermal energy of order
\begin{equation}
  E_{\mathrm{mode}}=kT.
  \label{eq:l9-equipartition}
\end{equation}

\begin{figure}[tbp]
  \centering
  \includegraphics[width=0.64\textwidth]{lecture_09_figure_02.png}
  \caption{The classical equipartition assignment for each electromagnetic
  normal mode of a thermal cavity.}
  \lectureframe{00:03:05}
\end{figure}

The cavity contains infinitely many modes, extending to arbitrarily high
frequency. If every one carries the nonzero energy
in Equation~\eqref{eq:l9-equipartition}, their sum diverges. The prediction is
not merely somewhat inaccurate in the ultraviolet; it is infinite. This is
the ultraviolet catastrophe.

Planck found the rule that cuts off the divergence:
\begin{equation}
  E_n=n h\nu,
  \qquad n=0,1,2,\ldots .
  \label{eq:l9-Planck-rule}
\end{equation}
At first, however, the rule was attached to oscillators in the cavity walls.
The radiation itself was still imagined as an ordinary classical wave.
Einstein took the more radical step and attributed the quanta to the
radiation. If \(h\nu\gg kT\), thermal agitation ordinarily cannot supply even
one quantum at that frequency. A classical fraction of a quantum is not
available; the most probable occupation is zero. High-frequency modes are
therefore suppressed.

\begin{figure}[tbp]
  \centering
  \includegraphics[width=0.64\textwidth]{lecture_09_figure_03.png}
  \caption{Planck's discrete energy rule, whose interpretation Einstein
  extended from material emitters to radiation itself.}
  \lectureframe{00:05:15}
\end{figure}

This did not make photons immediately welcome. A wave theory had explained
interference and diffraction too well to be discarded casually. The right
conclusion was not that waves were simply false, but that the classical wave
description was incomplete. Compton scattering eventually made the
particle-like transfer of energy and momentum especially difficult to evade.

\begin{classroomqa}
\audiencequestion{Why was the photoelectric effect not already decisive
evidence for photons?
\lecturetimestamp{00:10:09--00:10:19}}
\lecturerresponse{It was important evidence, and Einstein's explanation was
one of the earliest direct uses of light quanta. Historically, however, it
did not force everyone to abandon a classical field interacting with
quantized matter. Compton scattering displayed photon-like energy and
momentum conservation so transparently that the remaining classical
interpretations became much harder to sustain.}
\end{classroomqa}

\begin{classroomqa}
\audiencequestion{Where does Bohr's 1913 atomic model belong in this
transition?
\lecturetimestamp{00:10:20--00:11:13}}
\lecturerresponse{It is one of several sharp breaks rather than the final
theory. Bohr imposed discrete stationary orbits and quantum jumps on a
largely classical orbital picture. The complete conceptual transition came
later, especially through Heisenberg's matrix mechanics, Schr\"odinger's wave
mechanics, and Dirac's unifying operator language.}
\end{classroomqa}

Einstein had also recognized very early that atomic and radioactive events
carry an irreducibly statistical appearance. An excited atom can remain
undecayed for an unpredictable interval and then emit spontaneously. A
radioactive nucleus behaves similarly. One can predict distributions and
lifetimes with great precision without predicting the moment of an
individual event.

The old orbital models tried to preserve familiar mechanics by adding special
rules: only certain orbits were allowed, and radiation somehow occurred only
in jumps between them. Such models could organize data, but they did not
provide a consistent mechanics. Heisenberg's decisive move was to stop
describing unobservable electron orbits and build the theory from observable
transition quantities.

\begin{classroomqa}
\audiencequestion{Why did the old orbit models persist if they were already
in tension with classical mechanics?
\lecturetimestamp{00:12:50--00:13:22}}
\lecturerresponse{They retained a familiar picture and reproduced important
spectral regularities. That made them useful transitional models. Their
special rules were not a coherent dynamics, however, and increasingly
complicated atoms exposed the weakness. Heisenberg broke with the orbit
picture rather than adding another repair.}
\end{classroomqa}

\begin{classroomqa}
\audiencequestion{If every observation of a photon uses a quantized charged
particle, could quantizing only the matter be enough?
\lecturetimestamp{00:14:06--00:14:55}}
\lecturerresponse{Not in a consistent interaction theory. A quantum system
entangles with whatever it exchanges energy and momentum with. Leaving the
electromagnetic field classical while matter is quantum can imitate limited
experiments, but it cannot consistently describe emission, absorption,
fluctuations, and photon statistics together. Both sides of the interaction
must ultimately be treated quantum mechanically.}
\end{classroomqa}

\begin{classroomqa}
\audiencequestion{Why did Einstein later resist quantum mechanics after
recognizing spontaneous probability so early?
\lecturetimestamp{00:14:59--00:17:08}}
\lecturerresponse{He recognized the empirical need for probabilistic
transition laws, but he remained dissatisfied with probability as the final
description of an individual physical system. The issue was not whether the
statistical predictions worked; it was what they meant and whether a deeper
state of affairs had been omitted. His later objections were attempts to
show that the formalism, though successful, might not be complete.}
\end{classroomqa}

\begin{classroomqa}
\audiencequestion{Could probability merely reflect ignorance, as it does
when a classical coin is tossed?
\lecturetimestamp{00:17:09--00:18:19}}
\lecturerresponse{That possibility was precisely the source of the dispute.
Quantum mechanics assigns probabilities even when the state has been
specified as completely as the theory allows. The formalism predicts not
only uncertainty but interference between alternatives. Treating those
probabilities as ordinary ignorance does not reproduce the full structure
without adding a much more subtle underlying theory.}
\end{classroomqa}

The historical prelude fixes the direction of the argument that follows.
Quantum mechanics is the primary framework. Classical mechanics is not used
to derive it; at most, classical expressions suggest a Hamiltonian whose
quantum consequences we then work out. Finally we must recover classical
motion as an approximation to the quantum theory.

\section{The Free Hamiltonian as a Quantum Guess}

The abstract Schr\"odinger equation is
\begin{equation}
  i\hbar\frac{\partial}{\partial t}\ket{\psi(t)}
  =\hat H\ket{\psi(t)}.
  \label{eq:l9-abstract-Schrodinger}
\end{equation}
In the position representation, momentum is
\begin{equation}
  \hat p=-i\hbar\frac{\partial}{\partial x}.
  \label{eq:l9-momentum-operator}
\end{equation}
The two equations have a revealing parallel. A derivative with respect to
time is generated by the Hamiltonian; a derivative with respect to position
is generated by momentum. The sign in
Equation~\eqref{eq:l9-momentum-operator} is conventional, but once chosen it
must be used consistently.

\begin{figure}[tbp]
  \centering
  \includegraphics[width=0.64\textwidth]{lecture_09_figure_04.png}
  \caption{Time and space derivatives paired with the Hamiltonian and
  momentum operators.}
  \lectureframe{00:21:40}
\end{figure}

For a nonrelativistic free particle on a line, we guess
\begin{equation}
  \boxed{\hat H=\frac{\hat p^{\,2}}{2m}.}
  \label{eq:l9-free-H}
\end{equation}
This is a correspondence-guided guess, not a derivation of quantum mechanics
from classical mechanics. Free space has no preferred direction, so the
energy should be unchanged under \(p\mapsto-p\). The simplest even function
with the familiar low-speed behavior is \(p^2/(2m)\). The constant \(m\)
earns its interpretation as inertial mass through the resulting relation
between momentum and velocity.

\begin{figure}[tbp]
  \centering
  \includegraphics[width=0.58\textwidth]{lecture_09_figure_05.png}
  \caption{The nonrelativistic free-particle Hamiltonian chosen as an even
  function of momentum.}
  \lectureframe{00:24:00}
\end{figure}

The square of an operator means composition: \(\hat p^{\,2}\psi\) is
\(\hat p(\hat p\psi)\). Writing
\(\psi(x,t)=\braket{x}{\psi(t)}\), the probability of finding the particle
in an interval \(dx\) at a specified time is
\begin{equation}
  \Pr(x\in[x,x+dx];t)=|\psi(x,t)|^2\,dx.
  \label{eq:l9-position-probability}
\end{equation}
Position labels the possible outcomes of an observable. Time labels the
evolution of the state; in this elementary formulation it is not another
coordinate whose value is measured with probability \(|\psi|^2\).

Substituting Equations~\eqref{eq:l9-momentum-operator}
and~\eqref{eq:l9-free-H} into
Equation~\eqref{eq:l9-abstract-Schrodinger} gives
\begin{align}
  i\hbar\frac{\partial\psi}{\partial t}
  &=
  \frac{1}{2m}
  \left(-i\hbar\frac{\partial}{\partial x}\right)^2\psi\\
  &=
  -\frac{\hbar^2}{2m}\frac{\partial^2\psi}{\partial x^2}.
  \label{eq:l9-free-Schrodinger}
\end{align}
Thus the free-particle equation is
\begin{equation}
  \boxed{
  i\hbar\frac{\partial\psi(x,t)}{\partial t}
  =
  -\frac{\hbar^2}{2m}
  \frac{\partial^2\psi(x,t)}{\partial x^2}.}
  \label{eq:l9-free-Schrodinger-boxed}
\end{equation}

\begin{figure}[tbp]
  \centering
  \includegraphics[width=0.72\textwidth]{lecture_09_figure_06.png}
  \caption{The free Schr\"odinger equation obtained by applying the momentum
  operator twice.}
  \lectureframe{00:28:02}
\end{figure}

\section{Plane Waves and the Two Directions of Momentum}

We first seek a solution with definite momentum. Following the sign convention
used at the board, take
\begin{equation}
  \psi(x,t)=e^{ipx/\hbar}e^{i\omega t}.
  \label{eq:l9-plane-wave-ansatz}
\end{equation}

\begin{classroomqa}
\audiencequestion{Should the frequency \(\omega\) also appear in the spatial
exponent?
\lecturetimestamp{00:29:05--00:29:20}}
\lecturerresponse{No. The spatial exponent is fixed by the momentum
eigenvalue: \(-i\hbar\partial_x e^{ipx/\hbar}=p e^{ipx/\hbar}\).
The independent time factor is then determined by the Schr\"odinger equation.
If one writes the wave as \(e^{i(kx-\omega t)}\), the relation is
\(p=\hbar k\) and \(E=\hbar\omega\).}
\end{classroomqa}

The time derivative of Equation~\eqref{eq:l9-plane-wave-ansatz} is
\(\partial_t\psi=i\omega\psi\), while
\begin{equation}
  \frac{\partial^2\psi}{\partial x^2}
  =-\frac{p^2}{\hbar^2}\psi.
\end{equation}
The free Schr\"odinger equation therefore requires
\begin{equation}
  -\hbar\omega=\frac{p^2}{2m},
  \qquad
  \omega=-\frac{p^2}{2m\hbar}.
  \label{eq:l9-board-frequency}
\end{equation}
The negative value is a consequence of having written the time factor as
\(e^{+i\omega t}\). In the conventional positive-frequency notation,
\begin{equation}
  \boxed{
  \psi_p(x,t)
  =
  \exp\left[
    \frac{i}{\hbar}
    \left(px-\frac{p^2}{2m}t\right)
  \right],}
  \label{eq:l9-plane-wave}
\end{equation}
so \(E=p^2/(2m)\) and the time factor is \(e^{-iEt/\hbar}\).

\begin{figure}[tbp]
  \centering
  \includegraphics[width=0.74\textwidth]{lecture_09_figure_07.png}
  \caption{The plane-wave substitution fixes the time frequency through
  \(E=p^2/(2m)\).}
  \lectureframe{00:34:45}
\end{figure}

The same energy belongs to two momentum eigenstates:
\begin{equation}
  p_E=\pm\sqrt{2mE}.
  \label{eq:l9-two-momenta}
\end{equation}
The right-moving and left-moving waves differ in their spatial phases but
share the same time phase. Any superposition
\begin{equation}
  \psi(x,t)
  =
  \alpha\psi_p(x,t)+\beta\psi_{-p}(x,t)
  \label{eq:l9-degenerate-superposition}
\end{equation}
therefore remains in the same energy eigenspace. In a two-state,
box-normalized version of the problem, the probabilities for momenta
\(\pm p\) are \(|\alpha|^2\) and \(|\beta|^2\), while the probability of
the common energy is their sum.

\begin{figure}[tbp]
  \centering
  \includegraphics[width=0.76\textwidth]{lecture_09_figure_08.png}
  \caption{The \(+p\) and \(-p\) plane waves have a common energy because
  the Hamiltonian depends on \(p^2\).}
  \lectureframe{00:37:40}
\end{figure}

Ideal plane waves are not square-integrable on the whole line. They are
generalized, delta-normalized eigenstates. Normalizable particle states are
wave packets assembled from them.

\section{The General Free State in Momentum Space}

Linearity now gives the general free evolution as a Fourier superposition:
\begin{equation}
  \boxed{
  \psi(x,t)
  =
  \frac{1}{\sqrt{2\pi\hbar}}
  \int_{-\infty}^{\infty}dp\,
  e^{ipx/\hbar}
  e^{-ip^2t/(2m\hbar)}
  \widetilde\psi(p).}
  \label{eq:l9-general-free-solution}
\end{equation}
The coefficient \(\widetilde\psi(p)\) is the momentum-space amplitude at
\(t=0\). Every value of momentum evolves independently by the phase dictated
by its energy.

\begin{figure}[tbp]
  \centering
  \includegraphics[width=0.78\textwidth]{lecture_09_figure_09.png}
  \caption{The general free state as a continuous superposition of momentum
  eigenwaves.}
  \lectureframe{00:41:10}
\end{figure}

\begin{classroomqa}
\audiencequestion{In the integral, is \(p\) a particular measured momentum
or only a dummy variable?
\lecturetimestamp{00:41:55--00:42:13}}
\lecturerresponse{Inside the integral it is a dummy variable ranging over
all real momenta. A particular outcome may instead be denoted \(q\). Its
probability density is \(|\widetilde\psi(q)|^2\), independent of what symbol
was used for the integration variable.}
\end{classroomqa}

\begin{classroomqa}
\audiencequestion{What is held fixed while the momentum integral is
performed?
\lecturetimestamp{00:42:14--00:42:57}}
\lecturerresponse{The spacetime point \(x,t\) is fixed. We sum the complex
contribution made at that point by every momentum component. Constants such
as \(m\) and \(\hbar\) are fixed as well; only the integration variable runs.}
\end{classroomqa}

With the normalization of Equation~\eqref{eq:l9-general-free-solution},
\begin{equation}
  \Pr(p\in[q,q+dq])=
  |\widetilde\psi(q)|^2\,dq.
  \label{eq:l9-momentum-density}
\end{equation}

\begin{figure}[tbp]
  \centering
  \includegraphics[width=0.70\textwidth]{lecture_09_figure_10.png}
  \caption{The momentum-space Born rule:
  \(P(q)=\widetilde\psi^*(q)\widetilde\psi(q)\).}
  \lectureframe{00:44:24}
\end{figure}

\begin{classroomqa}
\audiencequestion{Does the symbol \(q\) in \(P(q)\) denote a position?
\lecturetimestamp{00:44:26--00:45:05}}
\lecturerresponse{No. Here \(q\) is merely a fresh name for a possible
momentum outcome. Position is still denoted \(x\). The change of letter
prevents confusion with the dummy momentum \(p\) already used inside the
integral.}
\end{classroomqa}

Energy is a coarser label than momentum because the map
\begin{equation}
  E=\frac{p^2}{2m}
  \label{eq:l9-energy-momentum}
\end{equation}
identifies \(p\) and \(-p\).

\begin{figure}[tbp]
  \centering
  \includegraphics[width=0.68\textwidth]{lecture_09_figure_11.png}
  \caption{For a given positive energy, the free Hamiltonian admits the two
  momentum roots \(p=\pm\sqrt{2mE}\).}
  \lectureframe{00:45:44}
\end{figure}

For discrete or box-normalized momentum states, the weight assigned to energy
\(E\) is the sum of the weights at its two roots. For a continuum
\emph{density} in energy, the change of measure must also be included:
\begin{align}
  \rho_E(E)
  &=
  \sum_{p_i:\,p_i^2/(2m)=E}
  \frac{|\widetilde\psi(p_i)|^2}
       {|dE/dp|_{p_i}}\\
  &=
  \frac{m}{\sqrt{2mE}}
  \left[
    |\widetilde\psi(\sqrt{2mE})|^2
    +
    |\widetilde\psi(-\sqrt{2mE})|^2
  \right],
  \qquad E>0.
  \label{eq:l9-energy-density}
\end{align}
The lecture's sum over the two momentum probabilities expresses the
essential degeneracy. Equation~\eqref{eq:l9-energy-density} adds the Jacobian
needed when probability is quoted per unit energy rather than for an energy
bin.

\begin{figure}[tbp]
  \centering
  \includegraphics[width=0.74\textwidth]{lecture_09_figure_12.png}
  \caption{The board collects both momentum amplitudes associated with one
  free-particle energy.}
  \lectureframe{00:48:20}
\end{figure}

\begin{classroomqa}
\audiencequestion{Should \(\hbar\) appear in the energy relation
\(E=p^2/(2m)\)?
\lecturetimestamp{00:47:58--00:48:24}}
\lecturerresponse{No. Momentum and energy retain the familiar
nonrelativistic relation. Planck's constant converts those quantities into
spatial and temporal phase rates: \(k=p/\hbar\) and
\(\omega=E/\hbar\). It belongs in the wave, not as an extra factor in the
dispersion relation.}
\end{classroomqa}

The physically admissible free states are precisely the suitable
superpositions of these modes. For ordinary one-particle wave packets,
\(\widetilde\psi\) belongs to \(L^2(\mathbb R)\), so the norm is finite.
Solutions that grow without bound and cannot represent finite probability
are not admitted as normalized states, although generalized eigenfunctions
remain useful as a basis.

\begin{classroomqa}
\audiencequestion{Does this claim of generality assume that every solution
is separable into a function of \(x\) times a function of \(t\)?
\lecturetimestamp{00:49:12--00:50:09}}
\lecturerresponse{No. Individual energy--momentum eigenmodes are separable,
but a general wave packet is their superposition and is usually not
separable. Completeness of the Fourier modes, not separability of the final
state, is what makes the integral general.}
\end{classroomqa}

\begin{classroomqa}
\audiencequestion{May \(\widetilde\psi(p)\) contain an arbitrary
momentum-dependent phase?
\lecturetimestamp{00:50:09--00:50:51}}
\lecturerresponse{Yes, and it matters. Its magnitude fixes the momentum
probability density, but its phase controls interference among momentum
components in position space. Two states with the same
\(|\widetilde\psi(p)|^2\) can therefore have very different spatial
wave packets.}
\end{classroomqa}

\section{Initial Data and Dispersive Evolution}

At \(t=0\), the time phase in
Equation~\eqref{eq:l9-general-free-solution} disappears:
\begin{equation}
  \psi(x,0)
  =
  \frac{1}{\sqrt{2\pi\hbar}}
  \int_{-\infty}^{\infty}dp\,
  e^{ipx/\hbar}\widetilde\psi(p).
  \label{eq:l9-Fourier-forward}
\end{equation}

\begin{figure}[tbp]
  \centering
  \includegraphics[width=0.72\textwidth]{lecture_09_figure_13.png}
  \caption{Setting \(t=0\) identifies the momentum superposition with the
  initial position-space wave function.}
  \lectureframe{00:52:02}
\end{figure}

Fourier inversion then gives
\begin{equation}
  \boxed{
  \widetilde\psi(p)
  =
  \frac{1}{\sqrt{2\pi\hbar}}
  \int_{-\infty}^{\infty}dx\,
  e^{-ipx/\hbar}\psi(x,0).}
  \label{eq:l9-Fourier-inverse}
\end{equation}
The factors of \(\hbar\) have been kept explicit so that \(px/\hbar\) is
dimensionless and the two representations have compatible normalization.

\begin{figure}[tbp]
  \centering
  \includegraphics[width=0.72\textwidth]{lecture_09_figure_14.png}
  \caption{The inverse Fourier transform recovers the momentum amplitude
  from the initial wave function.}
  \lectureframe{00:53:05}
\end{figure}

The initial-value problem is now constructive:
\begin{enumerate}
  \item begin with \(\psi(x,0)\);
  \item Fourier transform it to obtain \(\widetilde\psi(p,0)\);
  \item multiply each momentum component by
  \(e^{-ip^2t/(2m\hbar)}\);
  \item inverse transform to recover \(\psi(x,t)\).
\end{enumerate}
Equivalently,
\begin{equation}
  \widetilde\psi(p,t)
  =
  e^{-ip^2t/(2m\hbar)}
  \widetilde\psi(p,0).
  \label{eq:l9-momentum-time-evolution}
\end{equation}

\begin{figure}[tbp]
  \centering
  \includegraphics[width=0.74\textwidth]{lecture_09_figure_15.png}
  \caption{Free evolution in momentum space consists of an
  energy-dependent phase for each Fourier component.}
  \lectureframe{00:57:06}
\end{figure}

\begin{classroomqa}
\audiencequestion{Can the time dependence simply be included in a
time-dependent momentum wave function?
\lecturetimestamp{00:56:32--00:57:34}}
\lecturerresponse{Yes. Equation~\eqref{eq:l9-momentum-time-evolution} is
exactly that description. Its magnitude remains fixed for a free particle,
so momentum probabilities are conserved, while its phase changes at a rate
set by \(p^2/(2m\hbar)\).}
\end{classroomqa}

Because the phase is quadratic in \(p\), different components slip relative
to one another. A localized packet generally spreads. With
\(p=\hbar k\) and the conventional positive frequency
\begin{equation}
  \omega(k)=\frac{\hbar k^2}{2m},
\end{equation}
the phase velocity and group velocity are
\begin{equation}
  v_{\mathrm{ph}}=\frac{\omega}{k}=\frac{p}{2m},
  \qquad
  v_{\mathrm g}=\frac{d\omega}{dk}=\frac{p}{m}.
  \label{eq:l9-wave-velocities}
\end{equation}
The group velocity agrees with the classical velocity associated with the
central momentum of a narrow packet. The variation of \(v_{\mathrm g}\) over
the packet's momentum width is what produces dispersion.

\begin{classroomqa}
\audiencequestion{Is this the distinction between phase velocity and group
velocity?
\lecturetimestamp{00:57:35--00:58:17}}
\lecturerresponse{Yes. The crests of an individual Fourier component move
with \(v_{\mathrm{ph}}\), while the envelope of a narrow packet moves with
\(v_{\mathrm g}\). For the nonrelativistic quadratic dispersion relation
they are not equal, and components with different momenta have different
group velocities, so a general packet spreads.}
\end{classroomqa}

\begin{classroomqa}
\audiencequestion{Can the same free-particle equation be used for a photon?
\lecturetimestamp{01:00:05--01:00:55}}
\lecturerresponse{No. Equation~\eqref{eq:l9-free-H} is the
nonrelativistic massive-particle Hamiltonian. A photon obeys
\(E=c|p|\), not \(E=p^2/(2m)\), and has no rest-mass parameter of this
kind. Its propagation and polarization require a relativistic field
description.}
\end{classroomqa}

\section{Commutators and Poisson Brackets}

We now step back from the explicit free solution. For an operator \(\hat K\)
that may itself depend on time,
\begin{equation}
  \frac{d}{dt}\langle K\rangle
  =
  \frac{i}{\hbar}\langle[\hat H,\hat K]\rangle
  +
  \left\langle\frac{\partial\hat K}{\partial t}\right\rangle.
  \label{eq:l9-expectation-evolution}
\end{equation}
The lecture considers operators with no explicit time dependence, so the
last term vanishes. The remaining formula follows by differentiating
\(\bra{\psi}\hat K\ket{\psi}\) and applying the Schr\"odinger equation to
both the ket and the bra.

\begin{figure}[tbp]
  \centering
  \includegraphics[width=0.74\textwidth]{lecture_09_figure_16.png}
  \caption{The two differentiated state vectors combine into the
  expectation value of a commutator.}
  \lectureframe{01:04:50}
\end{figure}

\begin{classroomqa}
\audiencequestion{Does the sign depend on whether the commutator is written
as \([\hat H,\hat K]\) or \([\hat K,\hat H]\)?
\lecturetimestamp{01:03:11--01:05:17}}
\lecturerresponse{Yes, but the two standard forms are identical:
\[
  \frac{d}{dt}\langle K\rangle
  =
  \frac{i}{\hbar}\langle[\hat H,\hat K]\rangle
  =
  \frac{1}{i\hbar}\langle[\hat K,\hat H]\rangle.
\]
Reversing the order changes the sign of the commutator, and
\(1/i=-i\) changes it back.}
\end{classroomqa}

The classical law for a phase-space function \(K(x,p)\) is
\begin{equation}
  \dot K=\{K,H\}_{\mathrm{PB}},
  \label{eq:l9-classical-evolution}
\end{equation}
where
\begin{equation}
  \{A,B\}_{\mathrm{PB}}
  =
  \frac{\partial A}{\partial x}
  \frac{\partial B}{\partial p}
  -
  \frac{\partial A}{\partial p}
  \frac{\partial B}{\partial x}.
  \label{eq:l9-Poisson-definition}
\end{equation}
The order in Equation~\eqref{eq:l9-classical-evolution} follows directly
from the chain rule and Hamilton's equations:
\begin{align}
  \dot K
  &=
  \frac{\partial K}{\partial x}\dot x
  +
  \frac{\partial K}{\partial p}\dot p\\
  &=
  \frac{\partial K}{\partial x}
  \frac{\partial H}{\partial p}
  -
  \frac{\partial K}{\partial p}
  \frac{\partial H}{\partial x}.
\end{align}

\begin{figure}[tbp]
  \centering
  \includegraphics[width=0.74\textwidth]{lecture_09_figure_17.png}
  \caption{Hamilton's equations lead directly to
  \(\dot K=\{K,H\}_{\mathrm{PB}}\).}
  \lectureframe{01:09:31}
\end{figure}

The elementary classical identity
\begin{equation}
  \{f(x),p\}_{\mathrm{PB}}=f'(x)
  \label{eq:l9-classical-fp}
\end{equation}
has an exact quantum counterpart. Act with the commutator on an arbitrary
test wave function:
\begin{align}
  [f(\hat x),\hat p]\psi
  &=
  f(x)\left(-i\hbar\frac{d\psi}{dx}\right)
  +
  i\hbar\frac{d}{dx}\bigl(f(x)\psi(x)\bigr)\\
  &=i\hbar f'(x)\psi(x).
\end{align}
Because this holds for every wave function in the common operator domain,
\begin{equation}
  \boxed{[f(\hat x),\hat p]=i\hbar f'(\hat x).}
  \label{eq:l9-quantum-fp}
\end{equation}

\begin{figure}[tbp]
  \centering
  \includegraphics[width=0.74\textwidth]{lecture_09_figure_18.png}
  \caption{Applying the commutator to an arbitrary wave function exposes the
  derivative \(f'(\hat x)\).}
  \lectureframe{01:19:12}
\end{figure}

In the momentum representation the reciprocal identity is
\begin{equation}
  [\hat x,g(\hat p)]=i\hbar g'(\hat p),
  \label{eq:l9-xg}
\end{equation}
and the special case \(f(x)=x\) yields
\begin{equation}
  [\hat x,\hat p]=i\hbar.
  \label{eq:l9-canonical-commutator}
\end{equation}
These examples motivate the correspondence
\begin{equation}
  [\hat A,\hat B]
  \longleftrightarrow
  i\hbar\,\widehat{\{A,B\}_{\mathrm{PB}}}.
  \label{eq:l9-correspondence}
\end{equation}

\begin{figure}[tbp]
  \centering
  \includegraphics[width=0.70\textwidth]{lecture_09_figure_19.png}
  \caption{The canonical commutator and Poisson bracket display the basic
  quantum--classical correspondence.}
  \lectureframe{01:21:52}
\end{figure}

Both brackets are antisymmetric,
\begin{equation}
  [A,B]=-[B,A],
  \qquad
  \{A,B\}_{\mathrm{PB}}=-\{B,A\}_{\mathrm{PB}},
  \label{eq:l9-antisymmetry}
\end{equation}
and both act as derivations on products:
\begin{align}
  [AB,C]
  &=A[B,C]+[A,C]B,
  \label{eq:l9-commutator-Leibniz}\\
  \{AB,C\}_{\mathrm{PB}}
  &=A\{B,C\}_{\mathrm{PB}}
  +\{A,C\}_{\mathrm{PB}}B.
  \label{eq:l9-Poisson-Leibniz}
\end{align}
For classical functions the order of the ordinary factors is immaterial. For
operators it is not, which is why the placement of \(B\) in
Equation~\eqref{eq:l9-commutator-Leibniz} must be preserved.

\begin{figure}[tbp]
  \centering
  \includegraphics[width=0.76\textwidth]{lecture_09_figure_20.png}
  \caption{The product rules for Poisson brackets and commutators, with
  operator order retained in the quantum identity.}
  \lectureframe{01:29:55}
\end{figure}

\begin{classroomqa}
\audiencequestion{In the commutator product rule, should the second term be
\(B[A,C]\) or \([A,C]B\)?
\lecturetimestamp{01:29:31--01:30:55}}
\lecturerresponse{It is \([A,C]B\). Expanding
\(ABC-CAB\) and adding and subtracting \(ACB\) gives
\[
  ABC-ACB+ACB-CAB
  =A[B,C]+[A,C]B.
\]
The classroom correction matters because the factors need not commute.}
\end{classroomqa}

The correspondence in Equation~\eqref{eq:l9-correspondence} is exact for the
canonical elementary identities just derived, but it is not a universal
replacement rule for arbitrary functions. A classical monomial such as
\(x^3p\) has several inequivalent operator orderings. Symmetrization often
provides a natural choice, yet higher-order \(\hbar\) corrections remain in
general. The classical bracket is the leading correspondence limit, not a
license to reverse or rearrange arbitrary operator strings.

\section{A Potential and Ehrenfest's Theorem}

Add a position-dependent potential:
\begin{equation}
  \hat H=\frac{\hat p^{\,2}}{2m}+U(\hat x).
  \label{eq:l9-potential-H}
\end{equation}
In position space, \(U(\hat x)\) acts by multiplication, and the
Schr\"odinger equation becomes
\begin{equation}
  i\hbar\frac{\partial\psi(x,t)}{\partial t}
  =
  \left[
    -\frac{\hbar^2}{2m}\frac{\partial^2}{\partial x^2}
    +U(x)
  \right]\psi(x,t).
  \label{eq:l9-potential-Schrodinger}
\end{equation}
The associated force is
\begin{equation}
  F(x)=-U'(x).
  \label{eq:l9-force}
\end{equation}

\begin{figure}[tbp]
  \centering
  \includegraphics[width=0.76\textwidth]{lecture_09_figure_21.png}
  \caption{The Schr\"odinger equation for kinetic energy plus a
  position-dependent potential.}
  \lectureframe{01:34:48}
\end{figure}

\begin{classroomqa}
\audiencequestion{Does the momentum operator itself now depend on time?
\lecturetimestamp{01:35:01--01:35:40}}
\lecturerresponse{In the Schr\"odinger picture,
\(\hat p=-i\hbar\partial_x\) has no explicit time dependence. The state
changes, and therefore its momentum distribution and expectation value can
change. One may instead transfer time dependence to operators in the
Heisenberg picture, but that is a different representation of the same
physics.}
\end{classroomqa}

We now ask how the center of a wave packet moves. First set
\(\hat K=\hat x\) in
Equation~\eqref{eq:l9-expectation-evolution}. Since
\([U(\hat x),\hat x]=0\),
\begin{align}
  \frac{d}{dt}\langle x\rangle
  &=
  \frac{i}{\hbar}
  \left\langle
    \left[\frac{\hat p^{\,2}}{2m},\hat x\right]
  \right\rangle.
  \label{eq:l9-xdot-start}
\end{align}
The product rule and
\([\hat p,\hat x]=-i\hbar\) give
\begin{align}
  [\hat p^{\,2},\hat x]
  &=
  \hat p[\hat p,\hat x]+[\hat p,\hat x]\hat p\\
  &=-2i\hbar\hat p.
\end{align}
Consequently,
\begin{equation}
  \boxed{
  \frac{d}{dt}\langle x\rangle
  =
  \frac{\langle p\rangle}{m}.}
  \label{eq:l9-Ehrenfest-x}
\end{equation}

\begin{figure}[tbp]
  \centering
  \includegraphics[width=0.76\textwidth]{lecture_09_figure_22.png}
  \caption{The first Ehrenfest equation: average velocity equals average
  momentum divided by mass.}
  \lectureframe{01:42:18}
\end{figure}

Next choose \(\hat K=\hat p\). The kinetic term commutes with momentum, while
Equation~\eqref{eq:l9-quantum-fp} gives
\([U(\hat x),\hat p]=i\hbar U'(\hat x)\). Therefore
\begin{align}
  \frac{d}{dt}\langle p\rangle
  &=
  \frac{i}{\hbar}
  \left\langle
    [U(\hat x),\hat p]
  \right\rangle\\
  &=
  -\langle U'(\hat x)\rangle
  =
  \langle F(\hat x)\rangle.
  \label{eq:l9-Ehrenfest-p}
\end{align}

\begin{figure}[tbp]
  \centering
  \includegraphics[width=0.76\textwidth]{lecture_09_figure_23.png}
  \caption{The second Ehrenfest equation: the mean momentum changes by the
  expectation value of the force.}
  \lectureframe{01:45:39}
\end{figure}

Combining the two exact identities,
\begin{equation}
  m\frac{d^2}{dt^2}\langle x\rangle
  =
  -\langle U'(\hat x)\rangle.
  \label{eq:l9-exact-Ehrenfest}
\end{equation}
This resembles Newton's equation, but an important distinction remains:
\begin{equation}
  \langle U'(\hat x)\rangle
  \neq
  U'(\langle x\rangle)
  \quad\text{in general}.
  \label{eq:l9-force-distinction}
\end{equation}
If the packet is narrow compared with the length scale over which the force
changes, then
\begin{equation}
  \langle U'(\hat x)\rangle
  \approx U'(\langle x\rangle),
  \label{eq:l9-narrow-packet}
\end{equation}
and
\begin{equation}
  m\frac{d^2}{dt^2}\langle x\rangle
  \approx F(\langle x\rangle).
  \label{eq:l9-classical-limit}
\end{equation}
For potentials of at most quadratic degree, \(U'(x)\) is at most linear, so
the replacement in Equation~\eqref{eq:l9-narrow-packet} is exact for the
first moments even when the packet is broad. For a general nonlinear
potential, packet width and higher moments matter.

\begin{classroomqa}
\audiencequestion{Are all the apparent Poisson-bracket and commutator
analogies exact, and is \(\langle F\rangle\) the quantum version of the
resultant force?
\lecturetimestamp{01:46:36--01:47:35}}
\lecturerresponse{The elementary canonical relations and the two Ehrenfest
equations above are exact. A general replacement of every Poisson bracket by
a commutator is only a correspondence rule because operator ordering can
matter. In one dimension, \(\langle F(\hat x)\rangle\) is the probability
weighted average force acting across the packet. It equals the force at the
average position only under the additional conditions just stated.}
\end{classroomqa}

This is the promised order of explanation. Quantum mechanics supplies the
fundamental evolution. Classical motion emerges when the observables of
interest vary little across the state, so that averages close approximately
on themselves. We do not first assume a classical particle and then decorate
it with a wave; we derive the motion of a classical-looking center from the
quantum wave packet.

\section{What Has Been Established}

The historical crisis taught us not to confuse a successful approximation
with a fundamental mechanism. With that lesson in place, the free Hamiltonian
\(\hat p^{\,2}/(2m)\) led to the Schr\"odinger equation, plane waves, and the
dispersion relation \(E=p^2/(2m)\). Their Fourier superpositions gave every
normalizable free state and a direct procedure for evolving arbitrary
initial data.

The second half of the lecture exposed the algebra beneath the
quantum--classical relation. Commutators govern expectation values in the
same structural position that Poisson brackets govern classical
observables. For a particle in a potential, that parallel produces the exact
Ehrenfest equations
\begin{equation}
  \frac{d}{dt}\langle x\rangle=\frac{\langle p\rangle}{m},
  \qquad
  \frac{d}{dt}\langle p\rangle
  =-\langle U'(\hat x)\rangle.
\end{equation}
Only after these quantum statements are established do we take the
narrow-packet, smooth-force limit and recover Newton's equation for the
packet's center.
